% A worked example of tikzflora: a short multipage document that uses every
% layout. Compile with LuaLaTeX, twice so the table of contents fills in:
%
%   lualatex lorem && lualatex lorem      (or: latexmk lorem)
%
% Every \florasetup and \florathispage below is commented with what it
% does; the full list of keys is in README.md and on
% https://flora.alpacasoft.dev/guide.
\documentclass[11pt,twoside]{article}
% Two-sided pages with a wide outer margin: the flora grows there.
\usepackage[a4paper,inner=2cm,outer=5.6cm,top=4.2cm,bottom=3.5cm,
  headheight=14pt,headsep=18pt]{geometry}
\usepackage{lipsum}
\usepackage{fancyhdr}
\usepackage{tikzflora}

\pagestyle{fancy}
\fancyhf{}
\fancyhead[LE,RO]{\small\textsc{Florilegium Lorem}}
\fancyfoot[LE,RO]{\small\thepage}
\renewcommand\headrulewidth{0pt}

% The document's defaults. layout=split puts the flora in the outer margin:
% on the right of odd pages and the left of even ones. seed picks the
% garden; change it for a different one, or use seed=random. visitors=0.6
% invites a few bees, butterflies, birds and snails.
\florasetup{layout=split, seed=3, visitors=0.6}

\begin{document}

% --- Title page -----------------------------------------------------------
% A narrower, centred text block leaves room around the title. The title
% layout crowns the lines inside the floratitle environment.
\newgeometry{left=3.2cm,right=3.2cm,top=5cm,bottom=4cm}
\thispagestyle{empty}
\florathispage{layout=title}
\begin{floratitle}
\centering
{\Huge\scshape Florilegium Lorem\par}
\vspace{1.2em}
{\large A specimen in several layouts\par}
\end{floratitle}
\vspace{6cm}
\begin{center}
\begin{minipage}{0.7\textwidth}
\small\lipsum[1][1-5]
\end{minipage}
\end{center}
\clearpage

% --- Contents -------------------------------------------------------------
% surround grows a wreath on all four sides of the text block.
\thispagestyle{empty}
\florathispage{layout=surround}
\tableofcontents
\clearpage
\restoregeometry

% --- Body -----------------------------------------------------------------
% The body uses the document default, layout=split.
\section{Of gardens in the margin}
\lipsum[2-6]

\section{Of stems and their habits}
\lipsum[7-11]

% From here on: twice the lilies and umbels, no ghost blooms, a bird and a
% butterfly on most pages, and leaves in green inks. mix weights are 1 for
% the usual amount and 0 for none.
\florasetup{mix={lily=2, umbel=2, ghost=0, bird=1.5, butterfly=1.5},
  palette={leaves=green}}

\subsection{The habit of the climbing kind}
\lipsum[12-14]

% --- A centred page -------------------------------------------------------
% centered grows one plant in each side margin, on separate seeds. This
% page alone also paints its blooms in violet inks.
\clearpage
\newgeometry{left=4.6cm,right=4.6cm,top=4.2cm,bottom=3.5cm,headheight=14pt,headsep=18pt}
\florathispage{layout=centered, palette={blooms=violet, hero=violet}}
\section{Of symmetry, and its refusal}
\lipsum[15-17]
\clearpage
\restoregeometry

% --- Closing --------------------------------------------------------------
% A short last page: the flora grows into the space below the text, here
% at a lighter density.
\florathispage{density=0.6}
\section{Colophon}
Set with \textsf{tikzflora}, whose margins are grown afresh on every page.

It would not exist without two sources: the r/LaTeX thread
\emph{How do you make these flowers in LaTeX?}
(\texttt{old.reddit.com/r/LaTeX/comments/1xixznw}) and Kevin Raymond's
\emph{florilegium} (\texttt{github.com/kevinraymond/florilegium}), whose grown
margins, animation and visitors this package follows. The ink colours are
Pilot iroshizuku ramps measured by madoromi
(\texttt{madoromi.moonwit.ch}); the growth follows Wong, Zongker and Salesin,
\emph{Computer-Generated Floral Ornament} (SIGGRAPH 1998).

\end{document}
